% !TEX root = ../main.tex

\chapter{全文总结与未来展望}

\section{全文工作总结}

本文以 \textit{visualnav-transformer} 项目为基础，以 NoMaD 为高层核心算法，以云深处 Lite3 四足机器人为主要部署对象，围绕``基于动作扩散策略的足式机器人视觉目标导航与探索''展开了从算法到系统再到平台的完整研究。与只停留在离线仿真指标上的工作不同，本文始终把研究目标限定为算法、系统与平台能否形成可运行闭环，能否在真机实现部署。

围绕这一目标，本文主要完成了四方面工作。

第一，系统梳理了 NoMaD\cite{nomad2023} 的代码结构、训练目标与推理流程，明确了视觉编码器、距离预测头和扩散策略头之间的共享条件空间关系。在此基础上，本文进一步剖析了 NoMaD 采用扩散策略生成动作序列的核心机制，包括前向加噪、反向去噪的马尔可夫链采样过程，以及条件嵌入如何同时服务于距离估计与动作生成两个下游任务。这一梳理工作不仅为后续推理优化提供了清晰的模块边界，也为视觉编码器替换实验提供了可控的实验基线。

第二，围绕部署需求对高层推理进行了系统优化。具体而言，本文研究了四项关键技术：DDIM 推理加速通过缩减去噪步数在保持动作质量的同时大幅降低计算开销；CFG 引导增强通过调节无分类器引导强度改善目标条件响应；TTS 推理时搜索通过在采样阶段引入多候选评估与排序机制提升动作选择质量；视觉编码器升级通过将默认 EfficientNet-B0 替换为 DINOv2、ConvNeXt、ResNet 等候选编码器来比较视觉表征能力。本文为上述四项技术建立了统一的离线统计实验链路，使不同优化手段可以在同一评估框架下进行公平比较。这四项技术之间存在正交性：DDIM 主要影响推理速度，CFG 主要影响目标敏感度，TTS 主要影响动作多样性筛选，编码器升级主要影响视觉特征质量，因此它们可以组合使用以获得更优的综合性能。

第三，面向 Lite3 构建了完整的分层闭环系统。该系统由统一推理模块、有限状态机、导航主机、中层控制映射与平台接口五个层次组成。统一推理模块封装了 NoMaD 的前向推理与所有优化选项；有限状态机管理导航、恢复、停止等行为状态的切换逻辑；导航主机负责 topomap 节点选择与子目标规划；中层控制映射将高层航点动作转换为 Lite3 可执行的速度指令；平台接口则处理与 Lite3 底层 SDK 的通信协议。本文在 MuJoCo 仿真环境中对该系统进行了端到端验证，确认了从视觉输入到运动执行的完整数据通路的正确性。分层设计的核心优势在于各层职责明确、接口标准化，使得单层替换或升级不会影响其他层的正常工作。

第四，形成了面向 Jetson Orin 与 Lite3 真机平台的完整部署流程。该流程涵盖了模型导出与推理环境配置、相机标定与图像预处理对齐、topomap 采集与管理、安全控制与异常恢复等关键环节。同时，本文为真实实验的图像记录、轨迹可视化和任务组织设计了规范化接口，确保实验结果的可追溯性和可复现性。

从阶段性结论看，当前默认算法框架确定为 EfficientNet-B0 视觉编码器、DDIM-2 采样、默认 CFG=0 和 TTS-8 候选筛选。DDIM-2 是当前最适合进入闭环验证的推理配置，它在速度与质量之间取得了较好平衡；CFG 更适合作为可调部署参数而非默认固定选项，其最优值依赖于具体场景的目标分布特性；TTS-8 能够以较低额外代价改善候选轨迹选择，TTS-16 可作为算力预算充足时的增强项；四个视觉编码器训练权重已经完成离线统计与 MuJoCo 闭环接入验证，但``最终部署最优编码器''的结论仍需结合真实相机分布、边缘设备时延和真机安全约束进一步确认。

\section{主要贡献}

本文的主要贡献可归纳为以下三点：(1)建立了面向部署的 NoMaD 推理优化链路，将 DDIM、CFG、TTS 和视觉编码器替换纳入同一推理接口进行分析与控制。这一贡献的意义在于，现有工作通常只关注单一优化维度（如仅加速推理或仅替换编码器），缺乏将多种优化手段纳入统一框架进行系统评估的方法论。本文提出的统一链路使研究者能够在同一实验条件下量化不同优化组合的效果，为扩散策略在资源受限平台上的实用化部署提供了可操作的技术路径。(2)构建了适用于 Lite3 的分层闭环系统，用状态机与导航主机把高层 NoMaD 策略组织成可恢复、可切换、可扩展的任务链路。这一贡献填补了从``扩散策略算法验证''到``足式机器人系统集成''之间的工程空白。与直接将策略输出映射为电机指令的端到端方案不同，分层架构保留了各层的可解释性和可调试性，使系统在真实部署中具备故障定位与在线调整能力。(3)给出了 Orin 上的 Lite3 真机部署方案，并将图像记录、可视化与安全控制纳入统一工作流，提升了系统复现性。这一贡献为后续研究者提供了从算法到真机的完整参考实现，降低了足式视觉导航研究的实验门槛，同时规范化的实验管理流程也为跨平台迁移与多次重复实验奠定了基础。

\section{研究不足}

本文仍存在三方面不足，需要在后续工作中加以改进。

第一，视觉编码器实验虽然已经接入四个训练权重并完成离线统计与 MuJoCo 闭环比较，但闭环扫描中引入了 MuJoCo 域内的路线稳定项，因此该结果更适合证明四套权重均可进入系统闭环，而不能单独作为真实部署最优编码器的最终证据。最终选择仍需要结合 Jetson Orin 端真实推理时延、真实相机图像分布和真机安全实验来判断。

第二，真机实验的正式统计表和大规模重复实验尚未完整。虽然本文已在 Lite3 真机上完成了系统通路验证和初步导航任务，但受限于实验场地可用时间和设备调度约束，尚未积累足够数量的重复试验以支撑统计显著性分析。这导致第五章部分实验结果仍以定性描述和代表性案例为主，缺乏如成功率置信区间、轨迹偏差分布等定量统计指标。

第三，当前 topomap 仍主要依赖预先采集的视觉节点，尚未形成更强的在线更新与自维护能力。这一不足限制了系统在动态变化环境中的长期自主运行能力——当环境外观发生显著变化（如光照条件改变、障碍物移动）时，预采集的 topomap 节点可能失效，导致子目标匹配质量下降，进而影响导航成功率。

\section{未来展望}

后续工作可从以下方向继续推进：(1)完成视觉编码器在真实部署条件下的最终训练对比与选择。具体而言，需要在 Jetson Orin 与 Lite3 真机上分别测量 EfficientNet-B0、DINOv2-Small、ConvNeXt-Tiny、ResNet-50 等候选的推理延迟、图像分布鲁棒性、导航成功率和失败模式，从而形成兼顾性能与效率的最终选择。(2)补齐真机统计实验，建立离线指标与真机任务成功率之间的更明确对应关系。这要求设计标准化的实验场景集（包含不同难度等级的导航任务），在每个场景下执行足够次数的重复实验。(3)将当前统一推理模块与状态机扩展到 Tron1 等更多平台，验证分层结构的可迁移性。Tron1 作为双足人形机器人，其运动学约束与 Lite3 有本质区别，若分层系统仅需替换中层控制映射即可适配，则可有力证明本文架构设计的通用性与可扩展性。(4)引入语言目标与多模态条件，使系统从 image-goal 进一步扩展到更灵活的任务表达。结合视觉--语言模型将自然语言指令编码为目标条件，可以使机器人接受``去走廊尽头的红色门''等语义级导航指令，大幅提升系统在人机协作场景中的实用价值。


综合来看，本文已经把原本主要面向轮式平台验证的 NoMaD 高层视觉导航策略，扩展为一套面向 Lite3 四足平台的可分析、可仿真、可部署的研究链路：从算法层面的扩散策略推理优化，到系统层面的分层闭环架构设计，再到平台层面的真机部署与实验管理。
